% PATCH: N-05 (also closes the drafted part of R2-04)
% Target section: sections/results.tex (FigReal_Morphological_Transition_Energy.png caption +
%   energy paragraph, new tab:mode_steady_state_energy)
% Requirement: intake/pending/R02-01_metric_and_figure_audit.md §9.7 and §9.9 ("Informe 1") --
%   both explicitly marked "AUTORIZADA PARA EJECUTAR (2026-08-27)" by the author.
%
% §9.7 -- FigReal_Morphological_Transition_Energy.png (transition energy, previously rejected):
%   The figure was rewritten (experiments/_plotting/builders/energy_transition.py) to derive
%   power from |current_A| x 7.2V (fixed NiMH pack voltage, confirmed reliable by the hardware
%   team) integrated over the pre-calibrated electromechanical coupling window (2.92s for
%   Quadrupedal->Differential Mobile, 1.56s for the reverse), across 35 individually-detected
%   robot_cmd structural-command transitions -- not an inferred power-peak against no verified
%   event, which is why the prior version was rejected. Promoted with the same filename
%   (scripts/promote_figure.sh ... --force).
%
%   Caption (results.tex ~894):
%   Before: "Energetic footprint of physical morphological transitions. Integrated INA power
%     telemetry during the Quadrupedal -> Differential Mobile and Differential Mobile ->
%     Quadrupedal electromechanical coupling windows, across all hardware replicates."
%   After:  \revblue{"Morphological transition energy per direction, derived from
%     |I_current| x 7.2V (fixed NiMH pack voltage) integrated over the pre-calibrated
%     electromechanical coupling window (2.92s / 1.56s, matching Fig.~\ref{fig:FourPanelLatencies}),
%     across N=16 (Quadrupedal->Differential Mobile) and N=19 (Differential Mobile->Quadrupedal)
%     individually detected structural-command transitions. Error bars: +-1 SD; the
%     between-direction difference is not statistically significant (Welch's t~=1.18)."}
%
%   Paragraph (results.tex ~899): the old text claimed an "invariant energetic cost" for
%   Quad->Diff and "greater dispersion" for Diff->Quad, with a causal explanation (river-pebble
%   terrain, wheel slippage) written for the previously-rejected figure. The real numbers
%   contradict this directly -- SD is *higher* in Quad->Diff (2.93 J) than in Diff->Quad
%   (1.19 J), the opposite of the old claim -- so the causal narrative could not be kept; the
%   paragraph was rewritten around the real figures (5.06+-2.93 J, N=16; 4.13+-1.19 J, N=19;
%   Welch t~=1.18, not significant), with the observed extra dispersion in Quad->Diff stated
%   without attributing an unverified mechanical cause (per §9.7 step 3's explicit instruction).
%
% §9.9 -- R2-04 steady-state per-mode energy (author decision: table, not figure -- a 3-category
%   bar chart with one N=1 category communicates poorly, and a table is far cheaper to update
%   once real Omnidirectional data lands):
%   Added \label{tab:mode_steady_state_energy} immediately after the (now-rewritten) transition
%   energy paragraph, same subsection ("ENERGY CONSUMPTION"), same tabularx style as
%   tab:consolidated_physical_metrics. Quadrupedal (C, N=20, 1.94+-1.77 W) and Differential
%   Mobile (H, N=20, 1.36+-0.98 W) are real replicated measurements from
%   experiments/R2-04-mode-energy/; Omnidirectional (X) is a single-segment N=1 preliminary
%   estimate (5.71 W, no SD), marked as such in the caption and a table footnote -- per the
%   hardware team's 2026-08-27 confirmation that a dedicated Omnidirectional test battery is
%   coming, this row is explicitly NOT final. Appended one sentence to the transition-energy
%   paragraph connecting the table to the reviewer's explicit request, without claiming a
%   significant Quadrupedal-vs-Differential difference (Welch t~=1.26, not significant).
%
% PROGRESS.md: added row N-05 (this patch, status applied) for the §9.7 work; updated the
%   existing R2-04 row's Data source/Target section(s)/Patch file, status raised to `drafted`
%   (not `applied` -- the Omnidirectional row is a placeholder pending real replicated data, per
%   §9.9 step 3's explicit instruction not to close R2-04 while X stays N=1). When the real
%   Omnidirectional data lands: re-run experiments/R2-04-mode-energy/scripts/build_mode_energy.py,
%   update the table's X row and drop the preliminary footnote, and log that update as a new
%   `C-xx` correction row (not a new `N-xx`) per §9.9 step 4 -- only then does R2-04 move to
%   `applied`.
%
% Verification: scripts/validate_tex.sh and scripts/check_roundtrip.sh both pass after
%   reassemble.py. No LaTeX toolchain available locally to run compile_pdf.sh.
% Status: applied (N-05); R2-04 stays `drafted` pending real Omnidirectional data (see above).
